\documentclass[10pt,a4paper]{amsart}
\usepackage[margin=19mm]{geometry}
\usepackage{amsmath,amssymb,mathtools}
\usepackage[scr=rsfs,cal=euler]{mathalfa}
\usepackage{xfrac}
\usepackage[unicode,hidelinks,bookmarks=false]{hyperref}
\newcommand{\ie}{i.e.\ }
\newcommand{\cf}{cf.\ }
\newcommand{\resp}{respectively\ }
\newcommand{\Subsection}{\subsection}
\providecommand{\nobreakdash}{\nobreak\hbox{-}}
\setlength{\emergencystretch}{2em}
\multlinegap0pt
\usepackage{amssymb}
\usepackage{enumerate}
\newtheorem{lemma}{Lemma}
\newtheorem{theorem}{Theorem}
\title{Integral Numerical Radius and Operator Matrix Bounds}
\author{Shiva Sheybani, Hamid Reza Moradi,, Mohammad Sababheh}
\date{Source: arXiv:2602.14460v1 (CC BY 4.0)}
\begin{document}
\maketitle

\noindent\emph{Source and licence.} Integral Numerical Radius and Operator Matrix Bounds. Version arXiv:2602.14460v1 (CC BY 4.0), distributed by arXiv under the Creative Commons Attribution 4.0 International licence, http://creativecommons.org/licenses/by/4.0/, as recorded in the arXiv licence metadata for this exact version and shown on the abstract page https://arxiv.org/abs/2602.14460v1. Full licence notice: \texttt{licenses/arxiv-2602.14460-CC-BY-4.0.txt}.\par\smallskip

\noindent\textit{Editorial scope.} Counted unit: the article's theorem 10 (source0.tex lines 362--375). The article's complete original proof of it is delivered with it (source0.tex lines 376--426, one \texttt{proof} environment of 51 lines). No mathematics is written, repaired or re-derived: the statement and the proof are the article's own lines, in the article's own notation, and no retained line is substituted or re-flowed.  Retained uncounted as the source-local context the proof needs: lines 350--356.  Nothing was omitted from any retained span, so the package carries no omission record.  No retained line contains a citation, so the wrapper carries no bibliography and no bibliography entry.  No retained line contains a figure environment, an image-inclusion command or drawing code, so no image is delivered and the package declares no asset dependency.  Editorial formatting only, in addition: an AMS \texttt{amsart} preamble that restates the article's own declarations and macros used by the retained lines, namely the environments \texttt{lemma}, \texttt{theorem} on separate counters, together with the standard packages those lines need.  The retained environments print in this document as Lemma 1, Theorem 1 in order of appearance, whereas the article prints the same items with the numbering of its own full document.  The complete native source of the article as distributed in the arXiv e-print for this exact version is bundled unchanged as \texttt{sources/194708/source0.tex}.
\begin{lemma}\label{5}
Let $S,T\in \mathbb B\left( \mathbb H \right)$. Then
\[\left\| S+T \right\|\le 2\min \left\{ \int\limits_{0}^{1}{\omega \left( \left[ \begin{matrix}
   O & \left( 1-t \right)S+tT  \\
   t{{S}^{*}}+\left( 1-t \right){{T}^{*}} & O  \\
\end{matrix} \right] \right)dt},\int\limits_{0}^{1}{\left\| \left( 1-t \right)S+tT \right\|dt} \right\}.\]
\end{lemma}

\begin{theorem}\label{10}
Let $S,T\in \mathbb B\left( \mathbb H \right)$. Then
\[\begin{aligned}
  & \left\| S+T \right\|+\left| \int\limits_{0}^{1}{\left( \omega \left( \begin{matrix}
   O & \left( 1-t \right)S+tT  \\
   t{{S}^{*}}+\left( 1-t \right){{T}^{*}} & O  \\
\end{matrix} \right)-\left\| \left( 1-t \right)S+tT \right\| \right)dt} \right| \\ 
 & \le \int\limits_{0}^{1}{\left( \omega \left( \begin{matrix}
   O & \left( 1-t \right)S+tT  \\
   t{{S}^{*}}+\left( 1-t \right){{T}^{*}} & O  \\
\end{matrix} \right)+\left\| \left( 1-t \right)S+tT \right\| \right)dt} \\ 
 & \le \left\| S \right\|+\left\| T \right\|.
\end{aligned}\]
\end{theorem}

\begin{proof}
By Lemma \ref{5}, we have
\[\begin{aligned}
  & \left\| S+T \right\| \\ 
 & \le \int\limits_{0}^{1}{\omega \left( \left[ \begin{matrix}
   O & \left( 1-t \right)S+tT  \\
   t{{S}^{*}}+\left( 1-t \right){{T}^{*}} & O  \\
\end{matrix} \right] \right)dt}+\int\limits_{0}^{1}{\left\| \left( 1-t \right)S+tT \right\|dt} \\ 
 &\quad -\left| \int\limits_{0}^{1}{\omega \left( \left[ \begin{matrix}
   O & \left( 1-t \right)S+tT  \\
   t{{S}^{*}}+\left( 1-t \right){{T}^{*}} & O  \\
\end{matrix} \right] \right)dt}-\int\limits_{0}^{1}{\left\| \left( 1-t \right)S+tT \right\|dt} \right|.
\end{aligned}\]
So,
\begin{equation}\label{7}
\begin{aligned}
  & \left\| S+T \right\|+\left| \int\limits_{0}^{1}{\omega \left( \left[ \begin{matrix}
   O & \left( 1-t \right)S+tT  \\
   t{{S}^{*}}+\left( 1-t \right){{T}^{*}} & O  \\
\end{matrix} \right] \right)dt}-\int\limits_{0}^{1}{\left\| \left( 1-t \right)S+tT \right\|dt} \right| \\ 
 & \le \int\limits_{0}^{1}{\omega \left( \left[ \begin{matrix}
   O & \left( 1-t \right)S+tT  \\
   t{{S}^{*}}+\left( 1-t \right){{T}^{*}} & O  \\
\end{matrix} \right] \right)dt}+\int\limits_{0}^{1}{\left\| \left( 1-t \right)S+tT \right\|dt}. 
\end{aligned}
\end{equation}
We know that
\[\begin{aligned}
  & \omega \left( \left[ \begin{matrix}
   O & \left( 1-t \right)S+tT  \\
   t{{S}^{*}}+\left( 1-t \right){{T}^{*}} & O  \\
\end{matrix} \right] \right) \\ 
 & \le \frac{1}{2}\left( \left\| \left( 1-t \right)S+tT \right\|+\left\| t{{S}^{*}}+\left( 1-t \right){{T}^{*}} \right\| \right) \\ 
 & \le \frac{1}{2}\left( \left( 1-t \right)\left\| S \right\|+t\left\| T \right\|+t\left\| {{S}^{*}} \right\|+\left( 1-t \right)\left\| {{T}^{*}} \right\| \right) \\ 
 & =\frac{1}{2}\left( \left\| S \right\|+\left\| T \right\| \right).  
\end{aligned}\]
Thus,
\begin{equation}\label{8}
\int\limits_{0}^{1}{\omega \left( \left[ \begin{matrix}
   O & \left( 1-t \right)S+tT  \\
   t{{S}^{*}}+\left( 1-t \right){{T}^{*}} & O  \\
\end{matrix} \right] \right)dt}\le \frac{1}{2}\left( \left\| S \right\|+\left\| T \right\| \right).
\end{equation}
We also have
\[\left\| \left( 1-t \right)S+tT \right\|\le \left( 1-t \right)\left\| S \right\|+t\left\| T \right\|.\]
Hence
\begin{equation}\label{9}
\int\limits_{0}^{1}{\left\| \left( 1-t \right)S+tT \right\|dt}\le \int\limits_{0}^{1}{\left( \left( 1-t \right)\left\| S \right\|+t\left\| T \right\| \right)dt}=\frac{1}{2}\left( \left\| S \right\|+\left\| T \right\| \right).
\end{equation}
Now, combining two inequalities \eqref{8} and \eqref{9} together with \eqref{7} implies the desired result.
\end{proof}

\end{document}
